% P00 | Abstract | Abstract
% guide: dev_notes/paragraph_guides/P00_abstract.txt
We compare discrete-time and continuous-time models for pricing and hedging options on two problems: a one-year at-the-money European call on NVIDIA and an American put.
A scaled symmetric random walk converges to Brownian motion, with a Kolmogorov distance at the terminal time that falls like $N^{-1/2}$ with the exact constant \KSconst.
Path functionals such as the running maximum converge at the same slow rate, which is why discretely monitored barriers need a correction.
The CRR tree reaches the Black--Scholes price of \$\BSprice\ to within \$\CRRtenkErr\ at $N=10{,}000$ steps, whereas the Euler--Maruyama Monte Carlo price has a standard error of \$\MCbestSE\ at $5\cdot10^5$ paths.
A Leisen--Reimer tree with only 101 steps is accurate to \$\DeskLrErr, and a control variate reduces the Monte Carlo variance by a factor of \VRcv.
For the American put, a CRR tree, a finite-difference solver for the variational inequality and least-squares Monte Carlo agree on a value of about \TreeRichV\ for the baseline contract, and the early-exercise premium is \EEPatmPct\% of the European put at the money.
On the NVIDIA inputs the American put is worth \$\NvdaPutAm, \NvdaPutEEPpct\% more than the European put.
Daily delta hedging of the call leaves a hedging error with a standard deviation of \HedgeSDdailyPct\% of the premium, and a wrong volatility distorts the American exercise rule asymmetrically: exercising too early costs about twenty times as much as exercising too late in our worst cases.
